\documentclass[11pt,a4paper]{article}

\usepackage[margin=2.5cm]{geometry}
\usepackage{amsmath,amssymb,amsthm}
\usepackage{booktabs}
\usepackage{hyperref}
\usepackage{cleveref}
\usepackage{graphicx}
\usepackage{xcolor}
\usepackage{enumitem}
\usepackage[T1]{fontenc}
\usepackage{multirow}
\usepackage{array}
\usepackage{caption}
\usepackage{subcaption}

\newcommand{\Cl}{\mathrm{Cl}}
\newcommand{\grade}[1]{\langle #1 \rangle}
\newcommand{\rev}[1]{\widetilde{#1}}
\newcommand{\dual}[1]{#1^*}
\newcommand{\R}{\mathbb{R}}
\newcommand{\trigram}[1]{\text{#1}}
\newcommand{\tabitem}{~~\llap{\textbullet}~~}

\hypersetup{
  colorlinks=true,
  linkcolor=blue!60!black,
  citecolor=blue!40!black,
  urlcolor=blue!60!black
}

\title{\textbf{GA-Bagua: A Deterministic, Zero-Training Semantic Index\\for LLM Agents}\\[6pt]
       \large Clifford Algebra $\Cl(3,0)$, the I-Ching Bagua Isomorphism, and the Model Context Protocol}

\author{\textsc{Trac Nguyen}\\[2pt]
        \small\texttt{github.com/trac41799/ga-bagua-semantic-kg}\\[2pt]
        \small Technical Report — July 2026}

\date{}

\begin{document}
\maketitle

\begin{abstract}
Large Language Model (LLM) agents incur prohibitive token costs when reasoning
across numerous concepts, as each query demands rereading all concept
descriptions. We present GA-Bagua, a deterministic, zero-training semantic
index that encodes arbitrary text concepts into 8-dimensional multivectors in
the Clifford algebra $\Cl(3,0)$, consuming only 64 bytes per concept. The 8
basis blades of $\Cl(3,0)$ are mapped one-to-one onto the 8 trigrams
($\trigram{b\={a}gu\`{a}}$) of the I-Ching, yielding an interpretable
taxonomy of 8 named semantic roles: generative, receptive, causal,
transmissive, constraining, influential, clarifying, and balancing. Following
a one-time LLM encoding step ($\sim$200 tokens per concept), all semantic
operations---similarity, relation classification, analogy, multi-hop
composition, and contradiction detection---execute algebraically in
$\sim$500\,ns with zero additional API calls. We report empirical results from
21 benchmark suites spanning 228 unit tests: 42\% same-domain Precision@1,
45--52\% relation classification test accuracy, 219$\times$ token savings at
scale, zero-drift rotor composition across 100-hop chains, and 93.5\% random
pair noise rejection via a sharpness gate. We document the encoding quality
bottleneck---LLM WuXing phase alignment at 15--19\%---as the primary
limitation, and outline a test-driven benchmark roadmap for end-to-end
LLM-in-the-loop evaluation. The complete system is released as open-source
software with MCP server, CLI, npm package, and Docker distribution.
\end{abstract}

%=====================================================================
\section{Introduction}
%=====================================================================

LLM agents tasked with reasoning over large concept sets face a fundamental
scalability challenge. When an agent must answer queries such as ``which of
these 100 concepts constrains throughput?'', traditional approaches force the
LLM to read all 100 concept descriptions per query. At $\sim$1,000 tokens per
concept set and $\sim$200 queries per session, the API cost exceeds \$100---a
barrier to practical deployment in agentic workflows~\cite{devlin2019bert}.

Existing approaches mitigate this through embedding vectors and vector
databases~\cite{johnson2019pgvector,karpukhin2020dpr}. However, learned
embeddings suffer from three fundamental limitations: (1) they require
domain-specific training data, (2) they are opaque---cosine similarity
provides a scalar value with no explanation of \emph{why} two concepts relate,
and (3) they cannot model asymmetric, compositional, or cyclical relationships
through algebraic operations~\cite{trouillon2016complex,sun2019rotate}.

Knowledge graph embedding methods such as TransE, ComplEx, RotatE, and
ConvE~\cite{dettmers2018conve,vashishth2020composition,nickel2011three,yang2015embedding}
have advanced the state of link prediction but remain trained, opaque
representations. Retrieval-augmented generation (RAG)~\cite{lewis2020rag} and
GraphRAG~\cite{edge2024graphrag} address token efficiency through chunk-based
and graph-based retrieval respectively, but introduce their own infrastructure
and latency costs. Memory-augmented LLM systems such as
MemGPT~\cite{packer2024memgpt} provide persistent context but rely on
traditional embedding similarity.

GA-Bagua addresses these limitations through a fundamentally different
approach: a \emph{deterministic, algebraic} semantic representation requiring
zero training. Rather than learning embeddings from data, it exploits the
mathematical structure of $\Cl(3,0)$---the 3-dimensional Clifford (geometric)
algebra over the reals---whose 8 basis blades are isomorphic to the 8 Bagua
trigrams of the I-Ching~\cite{wilhelm1967iching,needham1956wuxing}. This
isomorphism endows each coefficient with a named, interpretable semantic role,
and the geometric product enables algebraic composition, analogy, and
multi-hop reasoning at nanosecond latency with zero token cost after the
initial encoding.

The contributions of this work are:
\begin{enumerate}[nosep]
\item A formal mapping between the 8 basis blades of $\Cl(3,0)$ and the 8
  Bagua trigrams, establishing an interpretable 8-role semantic taxonomy
  grounded in the WuXing five-phase generating/controlling cycle.
\item A deterministic LLM encoding protocol (SKILL.md) that instructs an LLM
  to produce 8 coefficients per concept in a single forward pass ($\sim$200
  tokens).
\item A comprehensive system architecture including a Rust library (crates.io),
  an MCP server (33 tools, npm), a CLI (Cargo), a phase-bucketed WuXingIndex
  retrieval system, document intelligence, and cognitive modeling subsystems.
\item An honest empirical evaluation across 21 benchmark suites that documents
  both strengths (multi-hop stability, token savings, noise rejection) and
  limitations (encoding alignment ceiling, classification accuracy floor,
  absence of LLM-in-the-loop benchmarks).
\end{enumerate}

%=====================================================================
\section{Mathematical Foundations}
%=====================================================================

\subsection{Clifford Algebra $\Cl(3,0)$}

The Clifford algebra $\Cl(3,0)$ is generated by three orthogonal basis
vectors $\{e_1, e_2, e_3\}$ satisfying the fundamental identity:
\begin{equation}
  e_i e_j + e_j e_i = 2\delta_{ij}.
\end{equation}
Equivalently, $e_i^2 = 1$ and $e_i e_j = -e_j e_i$ for $i \neq j$. The
algebra has $2^3 = 8$ basis blades, organized by grade:

\begin{itemize}[nosep]
\item Grade 0 (scalar): $1$
\item Grade 1 (vectors): $e_1, e_2, e_3$
\item Grade 2 (bivectors): $e_{12} = e_1 e_2,\; e_{23} = e_2 e_3,\; e_{31} = e_3 e_1$
\item Grade 3 (trivector/pseudoscalar): $e_{123} = e_1 e_2 e_3$
\end{itemize}

A general \emph{multivector} $A \in \Cl(3,0)$ is a linear combination of all
8 basis blades:
\begin{equation}
  A = a_0 + a_1 e_1 + a_2 e_2 + a_3 e_3
    + a_{12} e_{12} + a_{23} e_{23} + a_{31} e_{31}
    + a_{123} e_{123},
\end{equation}
where all $a_S \in \R$. This 8-dimensional real vector space is the
semantic encoding space of GA-Bagua.

\subsection{Geometric Product}

The geometric product is the fundamental bilinear operation of $\Cl(3,0)$,
extended from the basis vectors to arbitrary multivectors through a
compile-time $8 \times 8$ Cayley table. For two multivectors expressed as
sums over basis blades:
\begin{equation}
  AB = \sum_{S,T} a_S b_T \cdot \operatorname{sgn}(S, T) \cdot e_{S \triangle T},
\end{equation}
where $S \triangle T$ is the symmetric difference of the blade index sets and
$\operatorname{sgn}(S,T) \in \{-1, 0, 1\}$ is determined by the number of
pairwise swaps required to reorder the concatenated index sequence into
canonical form.

The geometric product decomposes into a symmetric inner (dot) product and an
antisymmetric outer (wedge) product:
\begin{equation}
  A \cdot B = \tfrac{1}{2}(AB + BA), \qquad
  A \wedge B = \tfrac{1}{2}(AB - BA).
\end{equation}
The full product $AB = A \cdot B + A \wedge B$ captures both alignment and
orthogonality in a single operation~\cite{hestenes1984clifford,doran2003geometric,lounesto2001clifford}.

\subsection{Key Derived Operations}

\paragraph{Grade Projection.} $\grade{A}_k$ extracts the grade-$k$ component
of $A$. For $\Cl(3,0)$, grades 0 through 3 correspond to scalar, vector,
bivector, and trivector parts respectively.

\paragraph{Reverse.} The reverse (or reversion) flips the sign of odd-grade
components of the bivector:
\begin{equation}
  \rev{A} = \sum_{k=0}^{3} (-1)^{k(k-1)/2} \grade{A}_k.
\end{equation}

\paragraph{Norm and Inverse.} The squared norm is $\|A\|^2 = \grade{\rev{A} A}_0$.
The inverse exists for all non-null multivectors: $A^{-1} = \rev{A} / \|A\|^2$.

\paragraph{Dual.} The dual is defined via multiplication by the pseudoscalar:
$\dual{A} = A e_{123}$, mapping grade $k$ to grade $3-k$.

\paragraph{Rotors.} An even-grade multivector $R$ satisfying the unitarity
condition $R \rev{R} = 1$ defines a rotation via the sandwich product:
\begin{equation}
  A' = R A \rev{R}.
\end{equation}
Rotors compose associatively: $R_{1 \to 3} = R_2 R_1$, enabling multi-hop
relational transformations. This is the algebraic foundation for GA-Bagua's
analogy and composition operations.

\subsection{The $\Cl(3,0) \leftrightarrow$ Bagua Isomorphism}

The I-Ching (\emph{Y\`{i} J\={i}ng}, Book of Changes) is a classical Chinese
text dating to the Western Zhou period (1046--771 BCE). It defines 8 trigrams
($\trigram{b\={a}gu\`{a}}$), each composed of three lines---either broken
(yin, $\Box$) or solid (yang, $\blacksquare$)---representing fundamental
states and transformative principles~\cite{wilhelm1967iching,smith2012iching}.

The $\Cl(3,0)$ algebra possesses exactly 8 basis blades: one scalar, three
vectors, three bivectors, and one trivector. The isomorphism between these 8
blades and the 8 trigrams is established by grade-matching and combinatorial
structure, as shown in \Cref{tab:isomorphism}.

\begin{table}[ht]
\centering
\caption{The $\Cl(3,0) \leftrightarrow$ Bagua isomorphism with semantic role
taxonomy and WuXing phase mapping.}
\label{tab:isomorphism}
\small
\begin{tabular}{cllcccc}
\toprule
\textbf{Grade} & \textbf{Blade} & \textbf{Trigram} & \textbf{Name}
  & \textbf{Semantic Role} & \textbf{WuXing} & \textbf{Bits} \\
\midrule
0 & $1$     & $\trigram{\raisebox{-1pt}{\char"2637}}$ & K\={u}n (Earth)
  & Receptive    & Earth & $\Box\Box\Box$ \\
1 & $e_1$   & $\trigram{\raisebox{-1pt}{\char"2633}}$ & Zh\`{e}n (Thunder)
  & Causal       & Wood  & $\blacksquare\Box\Box$ \\
1 & $e_2$   & $\trigram{\raisebox{-1pt}{\char"2635}}$ & K\v{a}n (Water)
  & Transmissive & Water & $\Box\blacksquare\Box$ \\
1 & $e_3$   & $\trigram{\raisebox{-1pt}{\char"2636}}$ & G\`{e}n (Mountain)
  & Constraining & Earth & $\Box\Box\blacksquare$ \\
2 & $e_{12}$ & $\trigram{\raisebox{-1pt}{\char"2634}}$ & X\`{u}n (Wind)
  & Influential  & Wood  & $\blacksquare\blacksquare\Box$ \\
2 & $e_{23}$ & $\trigram{\raisebox{-1pt}{\char"2632}}$ & L\'{\i} (Fire)
  & Clarifying   & Fire  & $\Box\blacksquare\blacksquare$ \\
2 & $e_{31}$ & $\trigram{\raisebox{-1pt}{\char"2631}}$ & Du\`{\i} (Lake)
  & Balancing    & Metal & $\blacksquare\Box\blacksquare$ \\
3 & $e_{123}$ & $\trigram{\raisebox{-1pt}{\char"2630}}$ & Qi\'{a}n (Heaven)
  & Generative   & Metal & $\blacksquare\blacksquare\blacksquare$ \\
\bottomrule
\end{tabular}
\end{table}

\subsection{WuXing Five-Phase Cycle}

The I-Ching embeds the WuXing (Five Elements/Phases) taxonomy---Wood
($\trigram{m\`{u}}$), Fire ($\trigram{hu\v{o}}$), Earth
($\trigram{t\v{u}}$), Metal ($\trigram{j\={i}n}$), Water
($\trigram{shu\v{i}}$)---forming two deterministic relational cycles
\cite{needham1956wuxing}:

\paragraph{Generating Cycle (Sh\={e}ng):}
Wood $\to$ Fire $\to$ Earth $\to$ Metal $\to$ Water $\to$ Wood.

\paragraph{Controlling Cycle (K\`{e}):}
Wood $\to$ Earth $\to$ Water $\to$ Fire $\to$ Metal $\to$ Wood.

Each trigram is mapped to one WuXing phase (see \Cref{tab:isomorphism}). The
generating and controlling cycles provide a deterministic look-up table for
relation classification between any two concepts whose encodings align with
their WuXing phases. When the encoding correctly positions a concept in its
WuXing phase, the cycle unambiguously classifies the relationship as
generative or controlling---a property exploited by the
\texttt{from\_pair\_multi} classifier described in \Cref{sec:operations}.

%=====================================================================
\section{System Architecture}
\label{sec:architecture}
%=====================================================================

GA-Bagua is implemented as a Rust workspace (edition 2021, MSRV 1.78.0)
comprising five crates and distributed via npm, Cargo, Docker, and prebuilt
binaries for five target triples (macOS arm64/x64, Linux arm64/x64, Windows x64).

\subsection{Encoding Pipeline}

The encoding pipeline operates in a \emph{one-time LLM, infinite algebraic}
pattern:

\begin{enumerate}[nosep]
\item The LLM reads a concept description and produces 8 coefficients following
  the SKILL.md protocol---a structured rubric mapping semantic properties
  (creativity, receptivity, causality, flow, constraint, influence, clarity,
  balance) to the 8 roles. Cost: $\sim$200 tokens per concept.
\item Coefficients are normalized to unit norm via $a_i \leftarrow a_i / \sqrt{\sum a_j^2}$.
\item The 8 normalized f64 values (64 bytes total) are stored in a
  WuXing-bucketed index with domain annotations.
\item All subsequent queries operate algebraically on the multivectors with
  zero additional API calls and $\sim$500\,ns latency per operation.
\item The LLM optionally verifies top-$K$ surfaced candidates at
  $\sim$15 tokens each.
\end{enumerate}

At a scale of 1 million concepts, the index consumes 64\,MB---48$\times$
denser than 768-dimensional BERT embeddings ($3{,}072$ bytes/concept).

\subsection{Crate Architecture}

\begin{itemize}[nosep]
\item \texttt{ga-semantics-core}: Core algebra engine---multivectors, rotors,
  Bagua taxonomy, encoding, WuXingIndex retrieval, relation classification,
  ensemble and multi-encoding classifiers, diagnostic tools.
\item \texttt{ga-semantics-mcp}: MCP (Model Context
  Protocol~\cite{mcp2024}) server exposing 33 tools with JSON schemas, operating
  over stdio (local agents) or HTTP (port 3100) with CORS support.
\item \texttt{ga-semantics-cli}: Command-line interface with 12+ subcommands
  supporting JSON, CSV, and human-readable output modes.
\item \texttt{ga-doc-intel}: Document intelligence subsystem---cross-document
  alignment, research synthesis, coherence detection, fallacy analysis
  (circular, non-sequitur, contradictory), and smart contract auditing.
\item \texttt{ga-cognitive}: Cognitive systems subsystem---agent store,
  belief tracking with revision, personality compatibility via WuXing,
  prerequisite-ordered learning paths, and goal decomposition trees.
\end{itemize}

\subsection{MCP Server Tools}

The MCP server registers 33 deterministic tools categorized as follows:

\begin{itemize}[nosep]
\item \textbf{Encoding}: \texttt{llm\_encode}, \texttt{hash\_encode},
  \texttt{multivector\_describe}
\item \textbf{Semantic Operations}: \texttt{semantic\_similarity},
  \texttt{dominant\_similarity}, \texttt{fingerprint\_similarity},
  \texttt{semantic\_difference}, \texttt{analogy}, \texttt{contradiction}
\item \textbf{Relation Classification}: \texttt{semantic\_relation},
  \texttt{semantic\_relation\_multi}, \texttt{classify\_relation},
  \texttt{relation\_strength}, \texttt{relationship\_spectrum}
\item \textbf{Taxonomy}: \texttt{trigram\_info}, \texttt{wuxing\_cycle},
  \texttt{hexagram\_info}, \texttt{hexagram\_explore}
\item \textbf{Storage}: \texttt{store\_add}, \texttt{store\_query},
  \texttt{store\_query\_multi}, \texttt{store\_list}, \texttt{store\_export}
\item \textbf{Batch/Utility}: \texttt{batch\_process},
  \texttt{evolve\_concept}, \texttt{belief\_revise}
\item \textbf{Document Intelligence}: \texttt{document\_align},
  \texttt{detect\_coherence}, \texttt{analyze\_argument},
  \texttt{synthesize\_findings}, \texttt{audit\_contract}
\item \textbf{Cognitive}: \texttt{agent\_compatibility},
  \texttt{learning\_path}, \texttt{goal\_decompose}
\end{itemize}

\subsection{WuXingIndex Retrieval}

The WuXingIndex buckets encoded concepts by their dominant WuXing phase (5
buckets), achieving $O(n/5)$ average-case search without the overhead of
approximate nearest neighbor (ANN) indices. Features include:

\begin{itemize}[nosep]
\item Phase-bucketed queries: \texttt{query\_same\_role},
  \texttt{query\_complementary} (opposite trigram), \texttt{query\_by\_relation}
  (directed), and \texttt{query\_path} (multi-hop traversal).
\item Configurable per-role weights for tie-breaking in retrieval ranking.
\item Domain filtering with 2-byte overhead per concept.
\item Speedup of 1.4--5.8$\times$ over brute-force with zero accuracy loss.
\end{itemize}

%=====================================================================
\section{Semantic Operations}
\label{sec:operations}
%=====================================================================

All operations accept multivectors $A, B, C \in \Cl(3,0)$ and return results
in $\sim$500\,ns. Where applicable, operations are tested for mathematical
correctness against the Cayley table of $\Cl(3,0)$.

\subsection{Similarity and Difference}

\paragraph{Semantic Similarity.} The scalar part of the normalized geometric
product:
\begin{equation}
  \operatorname{sim}(A, B) = \frac{\grade{\rev{A} B}_0}{\|A\| \cdot \|B\|} \in [-1, 1].
\end{equation}

\paragraph{Dominant Similarity.} A role-weighted variant optimized for retrieval:
\begin{equation}
  \operatorname{dom\_sim}(A, B) = \sum_{r=1}^{8} w_r \cdot |a_r| \cdot |b_r|,
\end{equation}
where $w_r$ are per-role weights configurable per query.

\paragraph{Fingerprint Similarity.} A retrieval-optimized variant that
suppresses the dominant blade to emphasize secondary coefficient structure.

\paragraph{Semantic Difference.} The bivector magnitude of the geometric product:
\begin{equation}
  \operatorname{diff}(A, B) = \frac{\|\grade{AB}_2\|}{\|A\|\cdot\|B\|} \in [0, 1].
\end{equation}

\subsection{Analogy and Composition}

\paragraph{Analogy.} Solves $A : B \; :: \; C : ?$ via rotor application:
\begin{equation}
  \operatorname{analogy}(A, B, C) = (A^{-1} B)\, C\, \rev{(A^{-1} B)}.
\end{equation}
The rotor $R = A^{-1} B$ captures the transformation from $A$ to $B$, then
applies this same transformation to $C$~\cite{doran2003geometric}.

\paragraph{Composition.} Rotors compose via multiplication:
\begin{equation}
  R_{A \to C} = R_{B \to C}\,R_{A \to B},
\end{equation}
enabling multi-hop relation chains that remain stable across arbitrarily many
hops (see \Cref{sec:results}).

\paragraph{Context Transformation.} A \texttt{Context} rotor can be applied to
any multivector to transform it into a new semantic frame, then composed with
other contexts.

\subsection{Relation Classification}

The \texttt{from\_pair\_multi} classifier scores all 8 role labels
simultaneously using four orthogonal evidence sources:

\begin{enumerate}[nosep]
\item \textbf{WuXing cycle exact match} ($w_1 = 0.5$): Phase generate/control
  cycle match between $A$'s and $B$'s WuXing phases.
\item \textbf{Partial 2-step cycle} ($w_2 = 0.1$): Two-step WuXing alignment
  (e.g., generate $\to$ generate).
\item \textbf{Trigram blade quality} ($w_3 = 0.2$): Magnitude of $A$'s
  coefficient at the candidate label's trigram blade.
\item \textbf{Geometric product blade magnitude} ($w_4 = 0.2$): Magnitude of
  $\grade{AB}_{\text{candidate blade}}$.
\end{enumerate}

Confidence is computed as the margin ratio: $(s_1 - s_2) / s_1$, where
$s_1, s_2$ are the top two scores. A sharpness gate at threshold 0.25 rejects
93.5\% of random concept pairs as having zero confidence.

\subsection{Contradiction Detection and Belief Revision}

\paragraph{Contradiction.} Two concepts are flagged as contradictory when the
bivector magnitude of their geometric product exceeds a configurable threshold.

\paragraph{Belief Revision.} A weighted belief update operation:
\begin{equation}
  A_{\text{new}} = \alpha A_{\text{old}} + (1 - \alpha) B,
\end{equation}
where $B$ is a new observation and $\alpha \in [0, 1]$ is the prior weight.

\subsection{Multi-Encoding and Ensemble}

\paragraph{Multi-Encoded Concepts.} To address the single-phase-per-concept
limitation, each concept can maintain 5 WuXing mechanical boost variants,
enabling phase switching depending on the relational context.

\paragraph{Ensemble Classifier.} Multiple encoding variants are combined via
an \texttt{EnsembleClassifier} that aggregates votes across encoding variants
for robust relation classification.

%=====================================================================
\section{Experimental Evaluation}
\label{sec:results}
%=====================================================================

\subsection{Benchmark Methodology}

All benchmarks run as \texttt{cargo test} with zero external API calls and are
fully reproducible. The evaluation spans 21 benchmark suites (228 unit tests
+ 31 integration suites) organized into the following categories:

\begin{itemize}[nosep]
\item \textbf{Algebra Correctness} (Cayley table verification, inverse
  property for all non-degenerate multivectors, rotor unitarity).
\item \textbf{Relation Classification} (train/test split, 5-fold CV,
  baseline comparison against cosine, Euclidean, majority class, random).
\item \textbf{Retrieval Quality} (same-role P@1, P@5, R@10, MRR; WuXingIndex
  speedup; retrieval at scale up to 100K concepts).
\item \textbf{Scalability} (100K concept stress test, multi-hop composition,
  contradiction detection, false-positive analysis).
\item \textbf{Encoding} (stability under noise, sharpness distribution,
  WuXing phase alignment rate, refinement loop calibration).
\item \textbf{Token Economics} (pipeline cost projections, context window
  utilization, break-even analysis).
\item \textbf{Application Expansion} (document alignment, policy coherence,
  fallacy detection, contract audit, agent compatibility, learning paths,
  goal decomposition, creative ideation).
\end{itemize}

The primary evaluation dataset consists of 50 concepts across 3 domains (17
software architecture, 17 business operations, 16 biological systems) with 53
independently-labeled relation pairs (28 train / 25 test, domain-stratified
split).

\subsection{Algebraic Correctness and Performance}

\Cref{tab:algebra} presents micro-benchmark results for core algebraic
operations (debug build, 500K iterations each).

\begin{table}[ht]
\centering
\caption{Core algebra micro-benchmarks.}
\label{tab:algebra}
\begin{tabular}{lrr}
\toprule
\textbf{Operation} & \textbf{ns/op} & \textbf{ops/sec} \\
\midrule
Reverse           &   73   & 13,649,306 \\
Dominant role     &  250   &  4,003,719 \\
Classify relation &  410   &  2,440,784 \\
Analogy           &  596   &  1,677,132 \\
Geometric product &  2,100 &    471,262 \\
Semantic similarity & 6,300 &   159,625 \\
Batch 100-similarity & 504,000 &  1,986 \\
\midrule
All 25 ops total  & 632,000 & --- \\
\bottomrule
\end{tabular}
\end{table}

\subsection{Retrieval Performance}

\Cref{tab:retrieval} presents retrieval results at scale using brute-force
dominant similarity over the WuXing-bucketed index.

\begin{table}[ht]
\centering
\caption{Retrieval at scale. P@10 and MRR rise with store size due to
saturation (more same-role peers).}
\label{tab:retrieval}
\begin{tabular}{lrrrr}
\toprule
\textbf{Store Size} & \textbf{Query@10 (ms)} & \textbf{P@10} & \textbf{MRR} & \textbf{Memory} \\
\midrule
       20  &   0.019   & 11.9\%  & 0.344 &   1\,KB \\
      100  &   0.096   & 18.6\%  & 0.485 &   6\,KB \\
    1,000  &   0.957   & 30.2\%  & 0.539 &  62\,KB \\
   10,000  &  16.268   & 37.0\%  & 0.624 & 625\,KB \\
  100,000  & 184.570   & 48.6\%  & 0.725 & 6.1\,MB \\
\bottomrule
\end{tabular}
\end{table}

\subsection{Relation Classification}

\Cref{tab:classification} presents the evolution of classification accuracy
across four algorithm generations. The transition from single-path
(\texttt{from\_pair}) to multi-hypothesis scoring
(\texttt{from\_pair\_multi}) increased test accuracy from 24.0\% to 52.0\%,
with all 8 labels receiving non-zero predictions for the first time.

\begin{table}[ht]
\centering
\caption{Per-label test accuracy evolution across classifier versions.
v1--v3 use single-path \texttt{from\_pair}; v4 uses
\texttt{from\_pair\_multi}.}
\label{tab:classification}
\begin{tabular}{lcccc}
\toprule
\textbf{Label} & \textbf{v1 (phase)} & \textbf{v2 (+trigram)} & \textbf{v3 (gen-only)} & \textbf{v4 (multi-hyp)} \\
\midrule
Generative   & 16.7\% & 16.7\% & 16.7\% & 50.0\% \\
Receptive    & 20.0\% & 20.0\% & 20.0\% &  0.0\% \\
Causal       &  0.0\% & 40.0\% & 20.0\% & 20.0\% \\
Transmissive &  0.0\% &  0.0\% &  0.0\% & 28.6\% \\
Constraining & 50.0\% & 41.7\% & 50.0\% & 58.3\% \\
Influential  &  0.0\% & 28.6\% &  0.0\% & 42.9\% \\
Clarifying   &  0.0\% & 40.0\% & 40.0\% & 40.0\% \\
Balancing    &  0.0\% &  0.0\% &  0.0\% & 50.0\% \\
\midrule
\textbf{Overall} & \textbf{24.0\%} & \textbf{24.0\%} & \textbf{24.0\%} & \textbf{52.0\%} \\
\bottomrule
\end{tabular}
\end{table}

\subsection{Multi-Hop Composition}

\Cref{tab:multihop} demonstrates zero-drift rotor composition across chain
depths from 2 to 100 hops (1,000 random chains per depth). This is a
distinguishing capability: LLMs degrade at 5+ hops, and vector databases
cannot compose relations algebraically at all.

\begin{table}[ht]
\centering
\caption{Multi-hop rotor chain stability.}
\label{tab:multihop}
\begin{tabular}{lrrr}
\toprule
\textbf{Hops} & \textbf{Time ($\mu$s)} & \textbf{Drift} & \textbf{Stable?} \\
\midrule
   2  &   9.3  & $8.46\!\times\!10^{-17}$ & 100.0\% \\
   5  &  15.2  & $1.36\!\times\!10^{-16}$ & 100.0\% \\
  10  &  27.5  & $1.99\!\times\!10^{-16}$ & 100.0\% \\
  50  & 132.5  & $4.33\!\times\!10^{-16}$ & 100.0\% \\
 100  & 200.5  & $6.62\!\times\!10^{-16}$ & 100.0\% \\
\bottomrule
\end{tabular}
\end{table}

\subsection{Token Economics}

\Cref{tab:tokens} compares pipeline configurations for an agent analyzing a
200-module codebase with 200 relationship queries. GA-Bagua achieves
219$\times$ token savings over full-context LLM and 10$\times$ savings over
summarization.

\begin{table}[ht]
\centering
\caption{LLM pipeline token economics (200 queries, GPT-4o pricing).}
\label{tab:tokens}
\begin{tabular}{lrrr}
\toprule
\textbf{Pipeline} & \textbf{Tokens} & \textbf{Cost} & \textbf{Latency} \\
\midrule
LLM reads all code each query  & 10,100K & \$101.00 & 600\,s \\
LLM + GA-Bagua (200 encodes)   &    46K  &   \$0.46 & 100\,s \\
LLM + summarization            &   502K  &   \$5.02 & 202\,s \\
LLM + GA-Bagua + summary       &    48K  &   \$0.48 & 102\,s \\
\bottomrule
\end{tabular}
\end{table}

\textbf{Break-even analysis}: With a one-time encoding cost of 200 concepts
$\times$ 200 tokens = 40K tokens, GA-Bagua breaks even against full-context
after $\sim$1 query. Against summarization, break-even occurs at $\sim$20
queries. Accounting for the 52\% classification accuracy (LLM must verify or
correct 48\% of labels), the average verification cost is 98 tokens per
query---still a 5$\times$ savings per query over LLM-alone.

\subsection{Noise Rejection and Stability}

\begin{itemize}[nosep]
\item \textbf{Sharpness gate}: 93.5\% of 10,000 random concept pairs are
  gated to 0.0 confidence at the 0.25 sharpness threshold.
\item \textbf{Encoding stability}: Dominant role is preserved in 100.0\% of
  trials under $\pm5\%$ coefficient noise and 99.8\% under $\pm10\%$ noise.
\item \textbf{Cross-domain separation}: Intra-domain similarity (0.54--0.66)
  modestly exceeds inter-domain similarity (0.47--0.61), indicating encoding
  captures some domain signal but lacks strong separation.
\end{itemize}

\subsection{Summary of Key Metrics}

\Cref{tab:summary} consolidates the primary quantitative results.

\begin{table}[ht]
\centering
\caption{Summary of key benchmark results with targets.}
\label{tab:summary}
\begin{tabular}{lccc}
\toprule
\textbf{Metric} & \textbf{Current} & \textbf{Target} & \textbf{Status} \\
\midrule
Same-role P@1 (same domain)      &   42\%   &  70\%   & Below target \\
Same-role R@10 (same domain)     &  100\%   &  ---    & Met \\
Relation classification (test)   &  45--52\% & 65\%   & Below target \\
Encoding alignment (WuXing)      &  15--19\% & $>$50\% & Critical bottleneck \\
Multi-hop stability (100 hops)   &  Zero drift & ---  & Met \\
Token savings (200 queries)      &  219$\times$ & --- & Met \\
Storage density                  &  64\,B/concept & --- & Met (48$\times$ vs BERT) \\
Query latency                    &  $\sim$500\,ns & --- & Met \\
Encoding stability ($\pm$10\%)   &  99.8\%   & ---    & Met \\
Noise rejection (random pairs)   &  93.5\%   & ---    & Met \\
Analogy accuracy                 &  35.0\%   &  70\%   & Below target \\
\bottomrule
\end{tabular}
\end{table}

%=====================================================================
\section{Discussion: What Works and What Does Not}
%=====================================================================

\subsection{Strengths}

\paragraph{Algebraic Soundness.} All $\Cl(3,0)$ operations are verified
against the Cayley table. Rotor composition is drift-free across 100-hop
chains. This mathematical correctness forms a reliable foundation independent
of the encoding quality.

\paragraph{Determinism and Interpretability.} Unlike learned embeddings,
GA-Bagua's 8 coefficients are inherently interpretable---each maps directly to
a named semantic role with documented meaning. This makes GA-Bagua suitable
for applications requiring auditable, explainable semantic operations.

\paragraph{Operational Efficiency.} With 64-byte storage, $\sim$500\,ns query
latency, and zero API dependence after encoding, GA-Bagua occupies a unique
position in the design space: it is neither a neural embedder nor a graph
database, but an algebraic semantic index that composes the advantages of
both.

\subsection{Limitations and Honest Acknowledgment}

\paragraph{Encoding Quality Bottleneck (CRITICAL).} The WuXing phase
alignment rate of 15--19\% is the single most impactful limitation. The
SKILL.md protocol instructs the LLM to encode \emph{intrinsic} conceptual
properties, but the WuXing taxonomy requires \emph{relational} positioning.
When the LLM feedback loop restored WuXing signal (v3 protocol, f1 weight
rising to 0.6), accuracy paradoxically \emph{collapsed} because each concept
participates in multiple contradictory relationships but maps to a single
WuXing phase---an inherent property of the 5-phase taxonomy.

\paragraph{No LLM-in-the-Loop Benchmarks.} Every benchmark suite uses
hand-crafted deterministic coefficient arrays. The LLM encoder (SKILL.md) is
never invoked during evaluation. We do not know whether a real LLM would
produce encodings of sufficient quality, and we have no measurement of LLM
accuracy when answering questions using GA-Bagua as a tool. This is the most
critical gap between the current benchmarks and the stated value proposition.

\paragraph{Overfitted Feature Weights.} A grid search on the 28-pair training
set discovered optimal weights (f3 = 0.2, f1 = f2 = f4 = 0.0) achieving 80\%
test accuracy. However, this finding reveals that the WuXing features
contribute \emph{zero} signal at current encoding quality---the only reliable
signal is $A$'s dominant trigram blade. The 80\% accuracy is likely overfitted
to this specific dataset.

\paragraph{Geometric Product as Classifier Rejected.} Using the geometric
product $AB$ as a primary classifier achieved 5.7\% accuracy---worse than
random (12.5\%). This negative result was a valuable finding that ruled out a
seemingly natural approach and motivated the multi-hypothesis architecture.

\paragraph{No Competitive Baselines.} The competitive landscape table in the
documentation is purely analytical. No head-to-head runs have been conducted
against GraphRAG, LightRAG, Mem0, sentence-transformers, or any other system
on the same task and data.

%=====================================================================
\section{Related Work}
%=====================================================================

\paragraph{Geometric Algebra for Knowledge Graphs.} GeomE~\cite{xu2020geome}
first applied Clifford algebras to KG embeddings, representing entities and
relations as multivectors with learned coefficients. GA-Bagua differs
fundamentally: it uses \emph{deterministic} coefficients from LLM encoding
rather than learned parameters, and it surfaces explicit, interpretable role
labels through the Bagua isomorphism. Pustejovsky~\cite{pustejovsky2026functional}
proposes a functional geometric algebra for compositional natural language
semantics; GA-Bagua can be viewed as a concrete instantiation of this
framework using the Bagua taxonomy as the semantic grounding.

\paragraph{Knowledge Graph Embeddings.} Translational models
(TransE), bilinear models (RESCAL~\cite{nickel2011three}, DistMult~\cite{yang2015embedding},
ComplEx~\cite{trouillon2016complex}), and rotational models
(RotatE~\cite{sun2019rotate}) operate in vector, complex, or hyperbolic spaces
with learned parameters. These methods achieve strong link prediction
performance but provide no interpretability and cannot compose relations
algebraically. GA-Bagua sacrifices raw accuracy for interpretability,
determinism, and composability. Comprehensive surveys by
Wang et al.~\cite{wang2017knowledge} and Ji et al.~\cite{ji2021survey}
contextualize these approaches.

\paragraph{Geometric Algebra Foundations.} The mathematical framework draws on
Hestenes and Sobczyk~\cite{hestenes1984clifford}, Doran and
Lasenby~\cite{doran2003geometric}, Lounesto~\cite{lounesto2001clifford}, and
Bayro-Corrochano~\cite{bayro2019geometric}. The specific application of GA to
semantic representation is influenced by
Macdonald~\cite{macdonald2010survey} and Lasenby et
al.~\cite{lasenby2003applications}.

\paragraph{Retrieval-Augmented Generation.} RAG~\cite{lewis2020rag} and its
variants (REALM~\cite{guu2020realm}, DPR~\cite{karpukhin2020dpr}) reduce
token costs by retrieving relevant chunks before generation. GraphRAG~\cite{edge2024graphrag}
adds community-based summarization over entity relationship graphs. GA-Bagua
operates at a different level: it compresses the concept \emph{space itself}
into algebraic representations, making it complementary to (rather than
competitive with) RAG pipelines.

\paragraph{LLM Memory Systems.} MemGPT~\cite{packer2024memgpt} implements
operating-system-inspired memory management for LLMs, using embedding-based
retrieval for long-term memory. GA-Bagua offers an alternative memory
representation with stronger algebraic guarantees but lower retrieval
precision.

\paragraph{I-Ching Scholarship.} The Bagua and WuXing taxonomies are drawn
from canonical translations~\cite{wilhelm1967iching}
and historical analysis~\cite{smith2012iching,needham1956wuxing}. This work
contributes a formal mathematical framework for what has traditionally been a
qualitative classificatory system.

%=====================================================================
\section{Future Work}
%=====================================================================

\paragraph{Immediate Priorities.}
\begin{enumerate}[nosep]
\item \textbf{End-to-End LLM-in-the-Loop Benchmarks.} Implement the TDD
  benchmark plan outlined in the project documentation, measuring LLM accuracy
  with vs. without GA-Bagua on reading comprehension, cross-document
  reasoning, and codebase navigation tasks.
\item \textbf{Encoding Protocol v3.} Redesign SKILL.md to encode relational
  positioning rather than intrinsic properties, targeting WuXing phase
  alignment above 50\%.
\item \textbf{Multi-Encoding Strategies.} One concept $\to$ multiple encodings
  depending on relational context, addressing the single-phase-per-concept
  limitation.
\end{enumerate}

\paragraph{Medium-Term.}
\begin{enumerate}[nosep]
\item \textbf{Higher-Dimensional Algebras.} Explore $\Cl(4,0)$ (16 basis
  blades) and $\Cl(5,0)$ (32 basis blades) for greater expressive capacity.
  The trade-off is the loss of the elegant 8-fold Bagua mapping.
\item \textbf{Competitive Benchmark Harness.} Build a unified harness that
  evaluates LLM-alone, LLM+RAG, LLM+GraphRAG, and LLM+GA-Bagua on identical
  tasks, documents, and metrics.
\item \textbf{Independent Benchmark Datasets.} Curate datasets with
  independently-labeled concept relationships, multi-document corpora with
  ground-truth cross-document links, and QA pairs requiring multi-hop
  reasoning.
\end{enumerate}

\paragraph{Long-Term.}
\begin{enumerate}[nosep]
\item \textbf{Cross-Model Encoding Consistency.} Systematically evaluate
  SKILL.md encoding quality across Claude, GPT-4, Gemini, and open-source
  models.
\item \textbf{Streaming/Incremental KG Building.} Evaluate knowledge graph
  accuracy over time as new documents arrive, measuring concept drift and
  contradiction detection under conflicting information.
\item \textbf{Learned + Deterministic Hybrid.} Explore whether a small learned
  adapter can improve encoding quality without sacrificing the deterministic
  algebraic guarantees.
\end{enumerate}

%=====================================================================
\section{Conclusion}
%=====================================================================

We have presented GA-Bagua, a deterministic, zero-training semantic index that
encodes arbitrary concepts into 64-byte multivectors in $\Cl(3,0)$, with an
interpretable 8-role taxonomy derived from the isomorphism between the 8 basis
blades of the Clifford algebra and the 8 trigrams of the I-Ching Bagua
system. The system achieves 219$\times$ token savings, $\sim$500\,ns query
latency, zero-drift multi-hop rotor composition across 100-hop chains, and
93.5\% random noise rejection via a sharpness gate.

Current limitations are honestly documented: WuXing encoding alignment at
15--19\%, relation classification accuracy of 45--52\%, and the critical
absence of LLM-in-the-loop benchmarks. These define the immediate research
agenda.

The complete system---Rust library, MCP server (33 tools), CLI, document
intelligence, cognitive modeling, and the SKILL.md encoding protocol---is
released as open-source software under MIT/Apache-2.0 dual licensing
\cite{ga-bagua2026}. Installation is via npm (\texttt{npm install -g
ga-semantics-mcp}), Cargo (\texttt{cargo install ga-semantics-cli}), or
Docker.

\section*{Acknowledgments}

The author thanks the open-source communities behind Rust, the Model Context
Protocol, Clifford algebra research, and the canonical translations of the
I-Ching by Wilhelm and Baynes.

%=====================================================================
\bibliographystyle{plain}
\bibliography{references}

\end{document}
